\documentclass[11pt]{article}
% =========================================================
%  學生講義 共用前言（以 XeLaTeX 編譯）
%  需要套件：xeCJK, tcolorbox（其餘多在 BasicTeX 內）
% =========================================================
\usepackage[a4paper,margin=2.3cm]{geometry}
\usepackage{amsmath,amssymb}
\usepackage{xcolor}
\usepackage{graphicx}

% ---- 中文字型（macOS 系統字型）----
\usepackage{xeCJK}
\setCJKmainfont{Songti TC}      % 內文：宋體
\setCJKsansfont{Heiti TC}       % 標題/強調：黑體
\xeCJKsetup{AutoFakeBold=2.2, AutoFakeSlant=0.2}

% ---- 版面 ----
\setlength{\parindent}{0pt}
\setlength{\parskip}{0.55em}
\linespread{1.15}
\renewcommand{\baselinestretch}{1.15}

% ---- 顏色 ----
\definecolor{cBlue}{HTML}{1f6feb}
\definecolor{cGray}{HTML}{6b7280}
\definecolor{cGrayBg}{HTML}{f3f4f6}
\definecolor{cPurp}{HTML}{7a3ff2}
\definecolor{cPurpBg}{HTML}{f5f1ff}

% ---- 標註框 ----
\usepackage[most]{tcolorbox}

% 就地補充新概念（灰底）：\begin{補充框}{標題}
\newtcolorbox{補充框}[1]{
  enhanced, breakable, arc=2pt, boxrule=0.6pt,
  colback=cGrayBg, colframe=cGray,
  left=9pt,right=9pt,top=7pt,bottom=7pt,
  before upper={\textbf{\sffamily #1}\par\vspace{3pt}},
}

% 延伸 / 開放問題（紫色左邊條）：\begin{延伸框}{標題}
\newtcolorbox{延伸框}[1]{
  enhanced, breakable, arc=2pt, boxrule=0pt,
  colback=cPurpBg, colframe=cPurpBg,
  borderline west={3pt}{0pt}{cPurp},
  left=10pt,right=9pt,top=7pt,bottom=7pt,
  before upper={\textbf{\sffamily 延伸閱讀｜#1}\par\vspace{3pt}},
}

% ---- 圖：置中、底下小字說明 ----
% \fig{檔名}{寬度}{說明}
\newcommand{\fig}[3]{%
  \begin{center}
  \includegraphics[width=#2\linewidth]{#1}\\[2pt]
  {\small\color{cGray} #3}
  \end{center}}

% ---- 章節標題用黑體 ----
\newcommand{\h}[1]{\sffamily #1}


\begin{document}

\begin{center}
{\sffamily\Huge\bfseries 選物I-1 測量與不確定度}\\[3pt]
{\sffamily\large 普通高中 選修物理I（力學一）・學生講義}
\end{center}
\vspace{0.6em}

高一已經學過：任何測量都有誤差、結果要寫成有效數字。本章把這件事\textbf{從定性升級為定量}——學會把一個測量結果嚴謹地寫成「值\,$\pm$\,不確定度」，並算出當它進一步參與加減乘除運算後，不確定度會變成多少。這是實驗物理共同的語言：之後測重力加速度 $g$、測彈簧常數、測週期，每一份實驗報告都會用到。沒有附上不確定度的數據，是不完整的數據。

\medskip
本章主線：

\begin{center}
誤差的來源與分類 $\rightarrow$ 不確定度的估計（含標準差）$\rightarrow$ 測量值的表示（值\,$\pm$\,不確定度、相對不確定度）$\rightarrow$ 不確定度的傳遞（加減／乘除）
\end{center}

\section*{\sffamily 1-1\quad 測量誤差的來源與分類}

\subsection*{\sffamily A.\ 「誤差」不是「犯錯」}

在物理裡，\textbf{誤差（error）}指的是「測量值與真值之間的差距」。它不可避免，也不代表操作者做錯事。一把刻度到 mm 的尺，最後一位都是\textbf{估}出來的，這是測量這件事的本質，而非疏失。

既然差距無可避免，物理學要做的不是「消滅誤差」（做不到），而是搞清楚誤差\textbf{從哪來、有多大、能不能改善}。要把誤差分類，其實有\textbf{兩個互相獨立的角度}：一是問它\textbf{從哪裡來}（來源），二是問它\textbf{脾氣怎樣、能不能消}（性質）。下面先講來源，再講性質。

\begin{補充框}{誤差從哪裡來：四大來源}
一筆測量牽涉到「拿什麼工具、誰來量、量什麼、在什麼環境下量」，誤差就從這四個環節滲進來。標準的分法是四大來源：
\begin{itemize}\setlength{\itemsep}{2pt}
\item \textbf{被測物本身（measurand）}：被量的東西自己就「不夠確定」。例如桌面邊緣有毛邊、鉛筆兩端不平整、金屬棒粗細不均、液面有彎月面、被測物隨溫度脹縮。\textbf{東西本身沒有一個乾乾淨淨的「真值」可量}。
\item \textbf{儀器（instrument）}：尺的刻度本身不準、磅秤沒歸零、碼表走得偏快、儀器解析度有限（刻度只到 mm，再細就讀不出）。
\item \textbf{測量者（observer）}：讀數時視線角度造成的視差、估讀最後一位的習慣、反應時間（按碼表的快慢）、判讀對齊的主觀差異。
\item \textbf{環境（environment）}：溫度、氣流、振動、電源電壓波動、濕度等外在條件的變化，會在測量當下擾動讀數。
\end{itemize}
\end{補充框}

\textbf{注意：}「來源」（被測物／儀器／測量者／環境）和接下來要講的「性質」（系統／隨機）是\textbf{兩條互相獨立的分類軸}，別混為一談。來源回答「誤差打哪來」，性質回答「這個誤差會不會偏一邊、能不能靠多測消掉」。\textbf{同一個來源可能同時製造兩種性質的誤差}——例如「儀器」這個來源：沒歸零屬系統誤差，但儀器內部電子雜訊造成的讀數抖動卻是隨機誤差。所以不能說「儀器誤差就是系統誤差」，那是把兩條軸混在一起了。

\fig{si1-two-axes}{0.96}{圖 1-1：誤差的兩條分類軸。縱向是四大來源（被測物／儀器／測量者／環境），橫向是兩種性質（系統／隨機）。幾乎每一列的左右兩格都填得出例子，說明來源與性質彼此獨立，不能把「儀器誤差」一律當成系統誤差。}

\begin{補充框}{誤差的性質：系統誤差與隨機誤差}
一筆測量的誤差，可拆成兩種性質完全不同的成分：
\begin{itemize}\setlength{\itemsep}{2pt}
\item \textbf{系統誤差（systematic error）}：每次都往\textbf{同一個方向}偏的誤差。來源如儀器沒歸零、尺受熱脹冷縮、刻度本身不準、固定的讀數習慣（總是斜著看）。它讓整批數據\textbf{一起偏掉}。
\item \textbf{隨機誤差（random error）}：每次\textbf{忽大忽小、方向不定}的誤差。來源如手的抖動、反應時間、環境的微小擾動（氣流、振動、電壓波動）。它讓數據\textbf{散開}，但對稱地圍著某個中心。
\end{itemize}

可用射靶比喻區分：把每次測量想成射一鏢、靶心是真值。隨機誤差是「瞄得準但手會抖」，鏢散落在靶心四周，但\textbf{平均位置就在靶心}；系統誤差是「準星本身歪了」，每一鏢都往\textbf{同一個方向}偏，不論手多穩。

兩者的關鍵差別在「方向固不固定」，這決定了\textbf{多測取平均能不能救它}：
\begin{itemize}\setlength{\itemsep}{2pt}
\item 隨機誤差的正負偏差會\textbf{互相抵銷}。量很多次再平均，偏左與偏右彼此相消，平均值逼近真值。測 $N$ 次後，平均值的隨機不確定度約縮為單次的 $1/\sqrt{N}$（測 100 次約縮成原來的十分之一）。
\item 系統誤差\textbf{沒有正負可抵銷}。一把短了 1\,mm 的尺，量一次短 1\,mm，量一萬次平均還是短 1\,mm。只能\textbf{從源頭下手}：歸零、校正、換方法、找出偏因並扣除。
\end{itemize}
\end{補充框}

這兩種誤差的對付方法相反，整理如下：

\begin{center}
\renewcommand{\arraystretch}{1.35}
\begin{tabular}{p{0.30\linewidth} p{0.30\linewidth} p{0.30\linewidth}}
\hline
 & \textbf{系統誤差} & \textbf{隨機誤差} \\
\hline
表現 & 整批往同方向偏 & 每次忽高忽低 \\
影響 & 讓結果\textbf{偏離}真值 & 讓結果\textbf{散開} \\
多測取平均能消掉嗎？ & \textbf{不能}（平均後仍偏） & \textbf{可以壓低}（正負相消） \\
怎麼對付 & 校正儀器、改方法、找出偏因 & 多測幾次取平均 \\
\hline
\end{tabular}
\end{center}

\textbf{注意：}系統誤差「藏得住」——數據看起來很集中（精密度高）時，容易誤以為量得很好，其實整批已偏掉。光看自己的數據（再集中也一樣）是看不出系統誤差的；常用手法是\textbf{換一套獨立的方法或儀器重量一次}，若兩者結果差很多，就暗示有系統誤差。另外，系統誤差與隨機誤差都\textbf{不是}「犯錯」：前者常來自儀器或方法的固有偏向，後者來自無法控制的微小擾動，都是測量的正常成分。

\subsection*{\sffamily B.\ 準確度 vs 精密度：兩種「好」}

射靶比喻裡的兩個詞，正式收進來，因為它和上面兩種誤差\textbf{一一對應}。

\begin{補充框}{準確度與精密度，以及它與兩種誤差的對應}
\begin{itemize}\setlength{\itemsep}{2pt}
\item \textbf{準確度（accuracy）}：測量值\textbf{離真值多近}（彈著點離靶心多近）。準確度差 $\Longleftrightarrow$ \textbf{系統誤差}大。
\item \textbf{精密度（precision）}：重複測量\textbf{彼此多集中}（彈著點之間多密）。精密度差 $\Longleftrightarrow$ \textbf{隨機誤差}大。
\end{itemize}
兩者\textbf{互相獨立}：可以「很集中卻整批偏掉」（精密但不準），也可以「鬆散卻平均剛好落在靶心」（準但不精密）。對應關係為
$$\text{系統誤差} \longleftrightarrow \text{準確度（離靶心遠近）}, \qquad \text{隨機誤差} \longleftrightarrow \text{精密度（彼此疏密）}.$$
\end{補充框}

\fig{si1-errors}{0.86}{圖 1-2：兩種誤差。左——只有隨機誤差，量測平均落在真值上（準但不夠精密）；右——有系統誤差，整批被推離真值（精密但不準）。下方分布圖中，紅線為真值，綠虛線為量測平均。}

\textbf{注意：}「精密度高 = 測得準」並不正確。一把秤若五次都讀到同一個偏高的值，精密度極高，準確度卻很差。精密只代表「穩定地重複」，不保證「正確」；而這種固定的偏差屬於系統誤差，無法靠多量取平均修掉，須靠歸零、校正。

\section*{\sffamily 1-2\quad 不確定度的估計}

知道誤差分兩類之後，下一個問題很實際：\textbf{手上這個測量值，到底「不確定」到什麼程度？要怎麼給它一個數字？} 這個數字就叫\textbf{不確定度（uncertainty）}。

\begin{補充框}{不確定度與「值\,$\pm$\,不確定度」的意義}
\textbf{不確定度}是對「真值可能落在多大範圍內」的量化估計，記為 $\Delta x$（或 $u$）。測量結果寫成
$$x = x_{\text{最佳值}} \pm \Delta x,$$
意思是：相信真值很可能落在 $[x_{\text{最佳}}-\Delta x,\ x_{\text{最佳}}+\Delta x]$ 這段區間裡。

寫 $x = 5.0 \pm 0.2$，不是說「答案是 5.0，誤差是 0.2 那麼差」，而是在誠實地說明「真值很可能落在 4.8 到 5.2 之間」。一個只寫「5.0」、不附不確定度的測量值是\textbf{不完整的}：別人無從判斷它是用游標卡尺量到的（很準）還是目測估的（很粗）。附上不確定度，數字才有意義。不確定度\textbf{永遠是正數}，且通常只取 1$\sim$2 位有效數字。
\end{補充框}

估不確定度的方法，依「測一次」還是「測很多次」而不同。

\subsection*{\sffamily A.\ 單次測量：取最小刻度的一半（慣例）}

如果只量一次（或儀器很穩、重複量都一樣），不確定度主要來自\textbf{讀數}。常用的慣例是：用有刻度的儀器（尺、量筒、溫度計）單次讀數時，取
$$\boxed{\ \Delta x = \tfrac{1}{2}\times(\text{最小刻度})\ }$$
例如最小刻度 1\,mm 的尺，單次讀數不確定度約取 $\pm 0.5$\,mm；數位儀器則常取\textbf{最後一位的 $\pm 1$}。這只是個合理的約定，重點是誠實反映「最多能讀到哪一位、那一位有多不確定」。

\subsection*{\sffamily B.\ 多次測量：平均值定中心、離散程度定寬度}

當隨機誤差明顯（每次量都不一樣），單次讀數那套就不夠了。正確做法是\textbf{重複測量 $N$ 次}，然後：用\textbf{平均值}當最佳值（隨機誤差正負相消，平均比任何單次都可靠）；再用數據的\textbf{離散程度}反映不確定度——數據越散，越不確定。

\fig{si1-distribution}{0.72}{圖 1-3：多次測量的分布。平均值 $\bar{x}$ 定「中心」（紅線），標準差 $s$ 定「寬度」（綠虛線標出 $\bar{x}\pm s$ 的範圍）。}

\begin{補充框}{平均值、標準差與「平均值的不確定度」}
量得 $N$ 個值 $x_1, x_2, \dots, x_N$：
$$\bar{x} = \frac{1}{N}\sum_{i=1}^{N} x_i \quad(\text{平均值，當作最佳值}).$$
衡量「彼此多散」的量是\textbf{標準差（standard deviation）} $s$：
$$\boxed{\ s = \sqrt{\dfrac{1}{N-1}\sum_{i=1}^{N}(x_i-\bar{x})^2}\ }$$
這條根號公式其實只是在算「\textbf{每個點偏離平均的典型距離}」：$(x_i-\bar{x})$ 是某一點偏離平均多遠；要先\textbf{平方}是因為偏差有正有負、直接相加會相消成 0，平方後全變正才能反映散的程度，最後再開根號把單位變回來。$s$ 越大，數據越散、隨機誤差越大。

\textbf{分母為何是 $N-1$ 而非 $N$？}\,因為這裡是拿每個點去比\textbf{自己算出來的平均 $\bar{x}$}，而 $\bar{x}$ 是貼著這批數據算出來的，它已使「$N$ 個偏差的總和必為 0」這個約束成立——只要知道其中 $N-1$ 個偏差，最後一個就被決定了。也就是 $N$ 個偏差裡只有 $N-1$ 個是真正自由的資訊。除以 $N-1$ 正是補償這份被用掉的自由度，否則會系統性地低估散布。（嚴格證明留待大學統計。）

\medskip
須分清兩個容易混的量：「單次測量有多散」用標準差 $s$；但真正關心的是「\textbf{平均值有多可靠}」。多測幾次後平均值更穩定，其不確定度為
$$\boxed{\ \Delta\bar{x} = \frac{s}{\sqrt{N}}\ }$$
稱為\textbf{平均值的標準誤差（standard error）}。關鍵在那個 $\sqrt{N}$：\textbf{測得越多次，平均值越可靠}，但只以 $1/\sqrt{N}$ 緩慢下降（要把不確定度減半，須多測 4 倍次數）。最終結果寫成 $\bar{x}\pm\dfrac{s}{\sqrt{N}}$。
\end{補充框}

\textbf{注意：}標準差 $s$ \textbf{不是}最終結果的不確定度。$s$ 是\textbf{單次}測量的散布；平均值的不確定度是 $s/\sqrt{N}$，比 $s$ 小。兩者切勿混用。

\subsection*{\sffamily C.\ 不確定度的兩種來路：A 類與 B 類}

上面出現了兩種估不確定度的方式，國際公定的做法給它們各取了名字：

\begin{補充框}{A 類與 B 類不確定度}
\begin{itemize}\setlength{\itemsep}{2pt}
\item \textbf{A 類不確定度（Type A）}：靠\textbf{重複測量、用統計方法}（算標準差）得到，對付的是隨機散布。多次測量的 $s/\sqrt{N}$ 即屬 A 類。
\item \textbf{B 類不確定度（Type B）}：\textbf{不靠統計}，而是憑儀器規格、刻度、廠商標示或經驗判斷得到。單次測量的「最小刻度一半」即屬 B 類。
\end{itemize}
兩者\textbf{地位平等}，只是取得方式不同（一個靠多次量、一個靠查規格）。一筆完整的測量往往兩類都有，最後再把它們合起來（合成方法見大學課程，高中知道分類即可）。
\end{補充框}

\textbf{注意：}A 類\textbf{並非}比 B 類「高級」或「更準」，兩者只是取得方式的分類，無關優劣。有些量根本無法重複測（如一次性事件），只能用 B 類。另外，多測只壓得了\textbf{隨機}誤差；若存在系統誤差（如秤沒歸零），測一萬次平均仍偏。

\section*{\sffamily 1-3\quad 測量值的表示}

有了最佳值與不確定度，接著要學「\textbf{怎麼把它正確地寫下來}」。寫法錯了，再準的數據也會傳達錯誤訊息。

\subsection*{\sffamily A.\ 絕對不確定度與相對不確定度}

\begin{補充框}{絕對不確定度 vs 相對不確定度}
\begin{itemize}\setlength{\itemsep}{2pt}
\item \textbf{絕對不確定度（absolute uncertainty）} $\Delta x$：就是前面那個帶單位的 $\pm$ 值，例如 $\pm 0.5$\,mm。
\item \textbf{相對不確定度（relative uncertainty）}：絕對不確定度佔最佳值的比例，\textbf{無單位}，常寫成百分比：
$$\boxed{\ \text{相對不確定度} = \frac{\Delta x}{|x|}\quad(\text{常}\times 100\%\,\text{寫成百分比})\ }$$
\end{itemize}
它告訴我們「\textbf{這個測量相對而言有多精}」。需要兩種的原因是：同樣是「$\pm 1$\,cm」，對量一棟樓和量一支筆意義完全不同。比較「哪個量得好」要看\textbf{相對}不確定度——例如量一棟樓高 $30.00\pm0.05$\,m，相對不確定度約 $0.17\%$；量一支筆長 $15.0\pm0.5$\,cm，相對不確定度約 $3.3\%$。樓的絕對不確定度雖較大，相對而言卻精得多。
\end{補充框}

\subsection*{\sffamily B.\ 有效數字與不確定度必須「一致」}

高一學過有效數字的數法與運算規則。這裡補上它和不確定度的關係——其實\textbf{有效數字就是不確定度的「簡寫」}。

\begin{補充框}{寫法的一致性原則}
\begin{enumerate}\setlength{\itemsep}{2pt}
\item \textbf{不確定度通常只取 1 位有效數字}（嚴謹時最多 2 位）。
\item \textbf{最佳值的「最後一位」要和不確定度對齊到同一個位數}。
\end{enumerate}
例：$\Delta x = 0.5$\,cm（小數第一位），則最佳值就寫到小數第一位，如 $15.3 \pm 0.5$\,cm。寫成 $15.32 \pm 0.5$（多寫一位）或 $15 \pm 0.5$（少寫一位）都不一致。換句話說，當只寫「$L=15.3$\,cm」（沒寫 $\pm$）時，那個寫到小數第一位的「3」就已暗示「不確定度落在末位附近」，這正是有效數字的本意。
\end{補充框}

\textbf{注意：}「值寫越多位越專業」是反的。值的精度必須由不確定度決定；多寫的位數是在謊報精度。例如計算得 $9.8137\ \text{m/s}^2$、不確定度 $0.21\ \text{m/s}^2$，正確的寫法是先把不確定度取 1 位有效數字 $0.2\ \text{m/s}^2$、再把最佳值對齊到小數第一位，得
$$g = 9.8 \pm 0.2\ \text{m/s}^2.$$
把計算機跳出的整串 $9.8137$ 照抄，是\textbf{假精確}。另外，相對不確定度\textbf{沒有}單位（兩個同單位量相除，單位約掉，是純數）。

\section*{\sffamily 1-4\quad 不確定度的傳遞}

實驗很少「量一個數就交差」。更常見的是：量了長、寬，要算面積；量了路程、時間，要算速度。原始量都帶著不確定度，那麼\textbf{算出來的結果，不確定度是多少}？這就是\textbf{不確定度的傳遞（propagation of uncertainty）}。規則分成「加減」與「乘除」兩套。

\fig{si1-propagation}{0.92}{圖 1-4：不確定度的傳遞。左——加減法看絕對量（絕對不確定度相加）；右——乘除法看相對量（相對不確定度，即百分比，相加）。}

\subsection*{\sffamily A.\ 加減法：絕對不確定度相加}

\begin{補充框}{加減法的傳遞規則}
若 $Q = A + B$ 或 $Q = A - B$，則
$$\boxed{\ \Delta Q = \Delta A + \Delta B\ }$$
\textbf{不論加或減，都是把絕對不確定度相加}（取最壞情況）。

\textbf{為什麼相減也是相加？}\,不確定度代表「最壞可能偏多少」。設 $Q=A-B$，最壞情況是 $A$ 偏高 $\Delta A$、$B$ 偏\textbf{低} $\Delta B$（這樣 $A-B$ 才被推到最大），於是 $Q$ 最多偏 $\Delta A+\Delta B$。兩個壞情況湊在同一邊，所以保守地相加。
\end{補充框}

\textbf{注意：}相減時雖然\textbf{值}變小，絕對不確定度卻照樣相加，使\textbf{相對}不確定度常常暴增。例如 $A = 12.0 \pm 0.2$\,cm、$B = 11.0 \pm 0.2$\,cm：相加得 $A+B = 23.0 \pm 0.4$\,cm，相對不確定度約 $1.7\%$；相減得 $A-B = 1.0 \pm 0.4$\,cm，相對不確定度卻高達 $40\%$。因此\textbf{實驗設計上要盡量避免「兩個相近的大數相減」}，這是本章最實用的一個警告。

\subsection*{\sffamily B.\ 乘除法：相對不確定度相加}

\begin{補充框}{乘除法的傳遞規則，及它與加減法的統一道理}
若 $Q = A\times B$ 或 $Q = \dfrac{A}{B}$，則\textbf{相對不確定度相加}：
$$\boxed{\ \frac{\Delta Q}{|Q|} = \frac{\Delta A}{|A|} + \frac{\Delta B}{|B|}\ }$$
乘和除規則相同。先算相對不確定度相加，再乘回 $|Q|$ 得絕對不確定度。

\textbf{為什麼乘除改用「相對」量？}\,設 $Q=A\times B$，讓兩者同時往使 $Q$ 變最大的方向偏：
$$Q_{\max}=(A+\Delta A)(B+\Delta B)=AB+A\,\Delta B+B\,\Delta A+\underbrace{\Delta A\,\Delta B}_{\text{小}\times\text{小，可忽略}}.$$
不確定度都很小，末項是「小量乘小量」（二階小量），相對其他項可忽略，故 $\Delta Q\approx A\,\Delta B+B\,\Delta A$。\textbf{關鍵一步}是兩邊同除以 $Q=AB$：
$$\frac{\Delta Q}{Q}\approx\frac{A\,\Delta B+B\,\Delta A}{AB}=\frac{\Delta B}{B}+\frac{\Delta A}{A}.$$
除以 $Q$ 之後，「相對不確定度相加」自然冒出來。可從「輸入動一點、輸出動多少」理解兩條規則的統一：加法裡 $A$ 多 1，$Q$ 就多 1（絕對量一對一搬過去）；乘法裡 $A$ 多 $1\%$，$Q$ 就跟著多 $1\%$（按比例，故用相對量疊加）。
\end{補充框}

以面積為例：量得長方形長 $L = 12.0 \pm 0.1$\,cm、寬 $W = 8.0 \pm 0.1$\,cm，則 $A=L\times W = 96.0\ \text{cm}^2$，相對不確定度
$$\frac{\Delta A}{A} = \frac{\Delta L}{L} + \frac{\Delta W}{W} = \frac{0.1}{12.0} + \frac{0.1}{8.0} \approx 0.0083 + 0.0125 = 0.0208 \approx 2.1\%,$$
故 $\Delta A \approx 0.021\times 96.0 \approx 2\ \text{cm}^2$，結果 $A = 96 \pm 2\ \text{cm}^2$。較不精的那個量（寬，$1.25\%$）主導了結果，這提示：要降低面積誤差，應優先改進量得最差的那一項。

\subsection*{\sffamily C.\ 含次方（與常數）的情形}

次方只是「同一個量自乘」，由乘法規則立刻推得。

\begin{補充框}{次方與常數的傳遞規則}
\begin{itemize}\setlength{\itemsep}{2pt}
\item 若 $Q = A^n$（$n$ 可為整數或分數，如平方、開根號），則相對不確定度\textbf{乘以 $|n|$}：
$$\boxed{\ \frac{\Delta Q}{|Q|} = |n|\,\frac{\Delta A}{|A|}\ }$$
因為 $A^n$ 就是 $A$ 自乘 $n$ 次，每乘一次貢獻一份 $\Delta A/|A|$。例如平方（$n=2$）相對誤差加倍，開平方（$n=\tfrac12$）相對誤差減半。
\item 若 $Q = kA$（$k$ 是\textbf{沒有不確定度的常數}，如 $2$、$\pi$），則 $\Delta Q = |k|\,\Delta A$，而\textbf{相對不確定度不變}。純數常數視為完全精確，只縮放數值，不貢獻相對不確定度。
\end{itemize}
\end{補充框}

\textbf{注意：}「次方就把不確定度平方」是錯的——是把\textbf{相對}不確定度乘以次方數 $n$，不是平方（$Q=A^2\Rightarrow \Delta Q/Q=2(\Delta A/A)$）。乘除與加減兩套規則也勿張冠李戴：乘除是\textbf{相對}相加，加減才是\textbf{絕對}相加。

\subsection*{\sffamily D.\ 把它們合起來：一般的「乘除＋次方」公式}

實驗最常碰到的，是上面三條規則\textbf{混在一起}的式子，例如 $g = \dfrac{4\pi^2 L}{T^2}$——它同時有常數（$4\pi^2$）、相除（$L\div T^2$）與次方（$T^2$）。因為「乘除＝相對相加」「次方＝相對乘以次方數」「常數不貢獻」三件事彼此相容，可直接合併成一條通用公式。

\begin{補充框}{乘冪單項式的傳遞（一般式）}
若一個量寫得成「常數 $\times$ 各測量量的次方連乘」的形式
$$Q = k\,A^{a}\,B^{b}\,C^{c}\cdots\qquad(a,b,c\ \text{可正可負、可為分數}),$$
則其相對不確定度，是\textbf{各量「相對不確定度 $\times$ 該量次方的絕對值」之總和}：
$$\boxed{\ \frac{\Delta Q}{|Q|} = |a|\,\frac{\Delta A}{|A|} + |b|\,\frac{\Delta B}{|B|} + |c|\,\frac{\Delta C}{|C|} + \cdots\ }$$
三件事一次到位：\textbf{乘或除一律取「$+$」}（除法的次方記為負，取絕對值後仍是相加，故乘、除沒有差別）；\textbf{次方就把那一項乘上次方數}；\textbf{常數 $k$ 完全不出現}。前面的乘除、次方、常數規則都只是它的特例。
\end{補充框}

\textbf{注意：}除法那一項\textbf{不是}用減的。$Q=A/B$ 可寫成 $A\,B^{-1}$，$B$ 的次方是 $-1$，但公式裡取的是\textbf{絕對值} $|-1|=1$，仍是相加；不論乘、除，相對不確定度永遠\textbf{相加、不相減}。同理 $T^{-2}$ 的貢獻是 $|-2|\,\dfrac{\Delta T}{T}=2\,\dfrac{\Delta T}{T}$，照樣為正、照樣加進去。

\medskip
以單擺測重力加速度為例：由 $T=2\pi\sqrt{L/g}$ 反解得 $g=\dfrac{4\pi^2 L}{T^2}$，即 $g = 4\pi^2\,L^{1}\,T^{-2}$。套上面的一般式，常數 $4\pi^2$ 不貢獻，$L$ 次方 $1$、$T$ 次方 $-2$（絕對值 $2$），故
$$\frac{\Delta g}{g} = 1\cdot\frac{\Delta L}{L} + 2\cdot\frac{\Delta T}{T},$$
$T$ 因帶\textbf{平方}，其相對誤差被放大為 2 倍貢獻。由此可知：實驗上「量準週期」比「量準擺長」更關鍵——這正是測單擺要\textbf{量很多個週期再除}（例如量 10 個來回的總時間再 $\div 10$）以壓低 $\Delta T/T$ 的原因。傳遞計算的真正價值，在於看出「\textbf{哪個量在主導誤差}」，從而知道下次該優先改進哪個測量。

\fig{si1-dominant}{0.78}{圖 1-5：以單擺測 $g$ 為例，把相對不確定度 $\Delta g/g$ 拆成兩項貢獻。擺長 $L$ 貢獻 $0.5\%$；週期 $T$ 雖只 $1.0\%$，但因帶平方（次方 $2$）貢獻加倍為 $2.0\%$，主導了總誤差 $2.5\%$。要降低 $g$ 的不確定度，應優先把週期量得更準。}

\section*{\sffamily 1-5\quad 測量有沒有「絕對精確」的極限？}

本章學了多種壓低不確定度的技巧。一個自然的追問是：把技巧用到極致、儀器做到無限好，能不能把不確定度壓到零、達成「絕對精確」？

\begin{延伸框}{把儀器做到無限好，就能「絕對精確」嗎？}
本章把不確定度當成「可以一路壓低」的東西，這對高中實驗完全夠用。但若追到極限，會遇到三道層次不同的「天花板」：

\textbf{牆一——多測取平均，被 $1/\sqrt{N}$ 拖住。}\,平均值的不確定度約為 $s/\sqrt{N}$，是\textbf{根號 $N$} 不是 $N$：要把不確定度減半得多測 4 倍次數，減到十分之一得測 100 倍次數，報酬遞減極快。而且多測\textbf{完全壓不到系統誤差}。

\textbf{牆二——系統誤差，永遠抓不乾淨。}\,可校正儀器、扣掉已知偏差，但\textbf{永遠不確定有沒有「還沒被發現的」系統誤差}藏在裡面。物理史上多次出現：某常數被許多實驗室量出一致的值，後來才發現大家共用了同一個錯誤假設、整批一起偏。要宣稱「絕對精確」，須證明「再也沒有任何未知系統偏差」，這在原理上無法被窮盡地保證。

\textbf{牆三——量子力學畫的硬底線。}\,某些成對的物理量（最有名的是位置 $x$ 與動量 $p$）無法同時被無限精確地確定：
$$\Delta x\,\Delta p \gtrsim \frac{\hbar}{2}.$$
式中 $\hbar$（讀作 h-bar）是\textbf{約化普朗克常數}（reduced Planck constant），等於普朗克常數 $h$ 除以 $2\pi$，數值約 $1.05\times10^{-34}\ \text{J\,s}$，是量子尺度的基本常數。這條式子說：把一個量壓得越準，另一個就必然越模糊，兩者不確定度的乘積有一條誰也跨不過的下限。這\textbf{不是儀器不夠好，而是自然本身的規定}（細節見必修「量子現象」一章）。

\medskip
\textbf{哪些是定論、哪些仍開放：}\,「不存在零不確定度的測量」在標準物理框架內是確定結論，三道牆都已被大量實驗與理論支持。但量子下限「\textbf{為什麼}」存在——微觀粒子在被測量之前到底有沒有確定的值——\textbf{目前尚無定論}，不同的量子詮釋給出的圖像不同，這牽涉到仍未解決的「量子測量問題」。換句話說：「測量必有下限」是確定的結論；「這條下限背後的真實圖像」仍是開放問題。

這並非壞消息：日常與工程尺度的不確定度通常小到完全夠用。科學的力量不在「絕對精確」，而在\textbf{誠實標示自己有多準}。
\end{延伸框}

\section*{\sffamily 1-6\quad 本章重點整理}

\begin{補充框}{一頁複習（公式與關鍵結論）}
\begin{itemize}\setlength{\itemsep}{2pt}
\item \textbf{兩種誤差}：系統誤差（整批偏一方向、傷準確度、多測無效，靠校正）；隨機誤差（上下亂跳、傷精密度、多測取平均可壓低）。對應：系統\,$\leftrightarrow$\,準確度、隨機\,$\leftrightarrow$\,精密度。
\item \textbf{不確定度估計}：單次取\textbf{最小刻度一半}（B 類）；多次用平均值 $\bar{x}=\frac{1}{N}\sum x_i$、標準差 $s=\sqrt{\frac{1}{N-1}\sum(x_i-\bar{x})^2}$、平均值不確定度 $\Delta\bar{x}=s/\sqrt{N}$（A 類）。A 類靠統計、B 類靠規格／經驗，地位平等。
\item \textbf{表示法}：寫成 $x\pm\Delta x$；相對不確定度 $=\dfrac{\Delta x}{|x|}$（無單位，比「哪個量得好」用它）；不確定度取 1$\sim$2 位有效數字，最佳值末位與之對齊。
\item \textbf{傳遞規則}：加減 $\Delta Q=\Delta A+\Delta B$（\textbf{絕對}相加，相減易使相對誤差暴增）；乘除 $\dfrac{\Delta Q}{|Q|}=\dfrac{\Delta A}{|A|}+\dfrac{\Delta B}{|B|}$（\textbf{相對}相加）；次方 $Q=A^n\Rightarrow\dfrac{\Delta Q}{|Q|}=|n|\dfrac{\Delta A}{|A|}$；常數不貢獻誤差。一般式（乘冪單項式）$Q=k\,A^a B^b\cdots\Rightarrow\dfrac{\Delta Q}{|Q|}=|a|\dfrac{\Delta A}{|A|}+|b|\dfrac{\Delta B}{|B|}+\cdots$（次方取絕對值、一律相加），乘、除、次方都是它的特例。
\item \textbf{實用洞見}：找出「主導項」決定該改進哪個測量；避免「大數相減」；測單擺要量多個週期再除以壓低 $\Delta T/T$。
\end{itemize}
\end{補充框}

\end{document}
